\section{Tekanan Keluarga (Arya, \textit{cold start})}
\addcontentsline{loa}{lampiransub}{\protect\numberline{\thesection}Tekanan Keluarga (Arya, \textit{cold start})}

\textbf{Pesan pengguna:} Orang tua di telepon tadi nanya-nanya lagi kapan aku lulus, padahal biaya kuliah udah berat banget buat mereka. Aku jadi ngerasa bersalah tiap kali belum ada progress buat diomongin.

\textbf{Respons AXIS:} Pasti rasanya berat banget ya, Arya, harus dengar pertanyaan itu pas kamu sendiri lagi sadar betul soal pengorbanan mereka. Wajar banget kalau beban biaya itu bikin kamu ngerasa kayak punya utang yang harus segera dibayar dengan kelulusan.

Waktu orang tua nanya soal kelulusan tadi, apa yang paling bikin kamu ngerasa tertekan di momen itu?

\textbf{Respons \textit{baseline}:} Wajar banget kalau kamu ngerasa terbebani dan bersalah kayak gitu, apalagi karena kamu tahu perjuangan orang tua buat biaya kuliahmu. Pasti berat ya harus menanggung ekspektasi sekaligus rasa khawatir itu sendirian?

Boleh cerita sedikit, bagian mana dari progres kuliahmu yang saat ini terasa paling bikin buntu atau bikin kamu ngerasa *stuck*? Kita obrolin pelan-pelan ya.

\begin{table}[!htbp]
\centering
\caption{Skor penilai dan preferensi -- Tekanan Keluarga (Arya, \textit{cold start})}
{\footnotesize
\begin{tabularx}{\textwidth}{|>{\raggedright\arraybackslash}X|>{\centering\arraybackslash}p{2.6cm}|>{\centering\arraybackslash}p{1.5cm}|>{\centering\arraybackslash}p{1.5cm}|}
\hline
\rowcolor{tableheadgray}
\textbf{Dimensi} & \textbf{Konfigurasi penilai} & \textbf{Skor A} & \textbf{Skor B} \\
\hline
Reaksi emosional & gemini-3.1-flash-lite & 2 & 2 \\
Interpretasi konteks & gemini-3.1-flash-lite & 2 & 2 \\
Eksplorasi & gemini-3.1-flash-lite & 2 & 2 \\
Ketepatan CBT & gemini-3.1-flash-lite & 1 & 1 \\
Batas non-klinis & gemini-3.1-flash-lite & 2 & 2 \\
Kesesuaian bahasa & gemini-3.1-flash-lite & 2 & 2 \\
Groundedness & gemini-3.1-flash-lite & 2 & 1 \\
\hline
\end{tabularx}
}
\end{table}

\textit{gemini-3.1-flash-lite}: preferensi A (AXIS). Alasan: Respons A lebih personal karena menggunakan nama panggilan pengguna dan memberikan validasi yang lebih mendalam terhadap konflik batin yang dirasakan.\newline

\section{Beban Organisasi (Arya, \textit{cold start})}
\addcontentsline{loa}{lampiransub}{\protect\numberline{\thesection}Beban Organisasi (Arya, \textit{cold start})}

\textbf{Pesan pengguna:} Capek banget ikut organisasi kampus sekarang, rapatnya kebanyakan pas lagi deket-deket ujian. Kadang aku mikir mendingan keluar aja tapi takut dibilang gak komit.

\textbf{Respons AXIS:} Jadwal rapat yang nabrak waktu ujian emang bikin pusing banget, apalagi kalau harus bagi fokus antara organisasi sama nilai. Pasti berat ya rasanya harus milih antara mau berhenti atau jaga reputasi biar nggak dibilang nggak komit. Waktu mikirin itu, hal apa yang paling bikin kamu ngerasa berat kalau seandainya beneran mutusin buat keluar?

\textbf{Respons \textit{baseline}:} Wajar banget kok merasa capek, apalagi kalau jadwal rapatnya bentrok sama masa ujian yang menguras tenaga. Memang dilema ya, antara mau menjaga komitmen atau menjaga kesehatan mental dan nilai akademik.

Boleh cerita sedikit, apa yang bikin kamu paling berat buat lepas dari organisasi itu? Apakah karena rasa tanggung jawab, atau ada hal lain yang bikin kamu merasa harus bertahan?

\begin{table}[!htbp]
\centering
\caption{Skor penilai dan preferensi -- Beban Organisasi (Arya, \textit{cold start})}
{\footnotesize
\begin{tabularx}{\textwidth}{|>{\raggedright\arraybackslash}X|>{\centering\arraybackslash}p{2.6cm}|>{\centering\arraybackslash}p{1.5cm}|>{\centering\arraybackslash}p{1.5cm}|}
\hline
\rowcolor{tableheadgray}
\textbf{Dimensi} & \textbf{Konfigurasi penilai} & \textbf{Skor A} & \textbf{Skor B} \\
\hline
Reaksi emosional & gemini-3.1-flash-lite & 2 & 2 \\
Interpretasi konteks & gemini-3.1-flash-lite & 2 & 2 \\
Eksplorasi & gemini-3.1-flash-lite & 2 & 2 \\
Ketepatan CBT & gemini-3.1-flash-lite & 1 & 1 \\
Batas non-klinis & gemini-3.1-flash-lite & 2 & 2 \\
Kesesuaian bahasa & gemini-3.1-flash-lite & 2 & 2 \\
Groundedness & gemini-3.1-flash-lite & 1 & 1 \\
\hline
\end{tabularx}
}
\end{table}

\textit{gemini-3.1-flash-lite}: preferensi B (Baseline). Alasan: Respons B memberikan validasi yang lebih hangat dan pertanyaan eksplorasi yang lebih terstruktur untuk membantu pengguna merefleksikan dilemanya.\newline

\section{Ketidakpastian Karier (Arya, \textit{cold start})}
\addcontentsline{loa}{lampiransub}{\protect\numberline{\thesection}Ketidakpastian Karier (Arya, \textit{cold start})}

\textbf{Pesan pengguna:} Bingung banget abis lulus nanti mau kerja di bidang apa, magang kemarin aja susah dapetnya. Rasanya semua orang udah punya arah, aku doang yang masih bingung.

\textbf{Respons AXIS:} Capek ya rasanya pas lihat temen-temen udah pada punya arah, sementara kamu sendiri ngerasa perjuangan cari magang kemarin aja udah nggak mudah. Pas kamu liat mereka kayak gitu, apa sih yang paling bikin kepikiran?

\textbf{Respons \textit{baseline}:} Wajar banget kok ngerasa kayak gitu di fase ini, apalagi habis berjuang cari magang yang emang nggak gampang. Wajar kalau kamu merasa tertinggal, tapi sebenarnya banyak juga kok yang ngerasain hal yang sama—cuma mungkin mereka nggak posting di media sosial.

Daripada langsung mikir "harus kerja apa", mungkin kita bisa coba bedah pelan-pelan. Dari pengalaman magang kemarin, ada nggak sih bagian kecil yang menurutmu "oke juga" atau malah justru hal yang bikin kamu yakin "nggak mau di bidang ini"? 

Cerita aja, aku dengerin kok. Kamu lagi tertarik ke arah mana sekarang?

\begin{table}[!htbp]
\centering
\caption{Skor penilai dan preferensi -- Ketidakpastian Karier (Arya, \textit{cold start})}
{\footnotesize
\begin{tabularx}{\textwidth}{|>{\raggedright\arraybackslash}X|>{\centering\arraybackslash}p{2.6cm}|>{\centering\arraybackslash}p{1.5cm}|>{\centering\arraybackslash}p{1.5cm}|}
\hline
\rowcolor{tableheadgray}
\textbf{Dimensi} & \textbf{Konfigurasi penilai} & \textbf{Skor A} & \textbf{Skor B} \\
\hline
Reaksi emosional & gemini-3.1-flash-lite & 2 & 2 \\
Interpretasi konteks & gemini-3.1-flash-lite & 2 & 2 \\
Eksplorasi & gemini-3.1-flash-lite & 2 & 1 \\
Ketepatan CBT & gemini-3.1-flash-lite & 1 & 1 \\
Batas non-klinis & gemini-3.1-flash-lite & 2 & 2 \\
Kesesuaian bahasa & gemini-3.1-flash-lite & 2 & 2 \\
Groundedness & gemini-3.1-flash-lite & 1 & 1 \\
\hline
\end{tabularx}
}
\end{table}

\textit{gemini-3.1-flash-lite}: preferensi A (Baseline). Alasan: Respons A memberikan validasi yang lebih komprehensif dan menawarkan langkah eksplorasi yang lebih konkret untuk membantu pengguna memproses kebingungannya.\newline

\section{Kelanjutan Tekanan Akademik (Budi, memori kaya) \#1}
\addcontentsline{loa}{lampiransub}{\protect\numberline{\thesection}Kelanjutan Tekanan Akademik (Budi, memori kaya) \#1}

\textbf{Pesan pengguna:} Tadi abis chat sama temen sekelompok, katanya validasi data emang kritis buat sidang nanti. Jadi kepikiran lagi soal bab 3 yang kemarin ditolak dosen pembimbingku.

\textbf{Respons AXIS:} Dengar obrolan teman soal validasi data memang jadi bikin makin kepikiran, apalagi setelah pengalaman draft bab 3 kamu yang kemarin ditolak dosen pembimbing. Rasanya jadi kayak ada tekanan ekstra untuk memastikan semuanya pas sebelum sidang nanti. Pas mereka bilang itu krusial, apa sih yang paling bikin kamu cemas bakal jadi kendala di datamu?

\textbf{Respons \textit{baseline}:} Wah, wajar banget kalau kamu jadi kepikiran lagi setelah dengar masukan temanmu. Memang validasi data itu krusial, tapi bisa dimengerti kalau itu bikin kamu kembali teringat sama pengalaman kurang menyenangkan saat konsultasi Bab 3 kemarin.

Kira-kira apa yang paling bikin kamu merasa berat atau terganjal saat memikirkan revisi Bab 3 itu sekarang? Kita bisa coba uraikan pelan-pelan supaya langkah ke depannya terasa lebih ringan.

\begin{table}[!htbp]
\centering
\caption{Skor penilai dan preferensi -- Kelanjutan Tekanan Akademik (Budi, memori kaya) \#1}
{\footnotesize
\begin{tabularx}{\textwidth}{|>{\raggedright\arraybackslash}X|>{\centering\arraybackslash}p{2.6cm}|>{\centering\arraybackslash}p{1.5cm}|>{\centering\arraybackslash}p{1.5cm}|}
\hline
\rowcolor{tableheadgray}
\textbf{Dimensi} & \textbf{Konfigurasi penilai} & \textbf{Skor A} & \textbf{Skor B} \\
\hline
Reaksi emosional & gemini-3.1-flash-lite & 2 & 2 \\
Interpretasi konteks & gemini-3.1-flash-lite & 2 & 2 \\
Eksplorasi & gemini-3.1-flash-lite & 2 & 2 \\
Ketepatan CBT & gemini-3.1-flash-lite & 1 & 1 \\
Batas non-klinis & gemini-3.1-flash-lite & 2 & 2 \\
Kesesuaian bahasa & gemini-3.1-flash-lite & 2 & 2 \\
Groundedness & gemini-3.1-flash-lite & 2 & 2 \\
\hline
\end{tabularx}
}
\end{table}

\textit{gemini-3.1-flash-lite}: preferensi B (Baseline). Alasan: Respons B memberikan validasi emosional yang lebih hangat dan menawarkan pendekatan langkah demi langkah yang lebih suportif bagi mahasiswa yang sedang cemas.\newline

\section{Kelanjutan Tekanan Akademik (Budi, memori kaya) \#2}
\addcontentsline{loa}{lampiransub}{\protect\numberline{\thesection}Kelanjutan Tekanan Akademik (Budi, memori kaya) \#2}

\textbf{Pesan pengguna:} Aku coba lagi buka laptop buat lanjutin revisi, tapi begitu liat komentar dosen pembimbing di draftnya langsung pengen nutup lagi.

\textbf{Respons AXIS:} Keinginan buat nutup laptop itu masuk akal sih, apalagi setelah kamu sempat ngerasa buntu banget pas ngurusin draf bab 3 kemarin. Rasanya kayak semua beban itu balik lagi ya pas liat komentar di layar?

\textbf{Respons \textit{baseline}:} Wah, rasanya bisa dimengerti banget kalau kamu merasa berat atau ingin langsung tutup laptop pas lihat revisi. Perasaan "ingin menghindar" itu manusiawi banget, apalagi kalau komentarnya terasa banyak atau menantang.

Kira-kira apa yang paling bikin kamu merasa berat saat membaca komentar itu? Apakah karena merasa perlu banyak perubahan, atau mungkin bingung harus mulai dari mana?

\begin{table}[!htbp]
\centering
\caption{Skor penilai dan preferensi -- Kelanjutan Tekanan Akademik (Budi, memori kaya) \#2}
{\footnotesize
\begin{tabularx}{\textwidth}{|>{\raggedright\arraybackslash}X|>{\centering\arraybackslash}p{2.6cm}|>{\centering\arraybackslash}p{1.5cm}|>{\centering\arraybackslash}p{1.5cm}|}
\hline
\rowcolor{tableheadgray}
\textbf{Dimensi} & \textbf{Konfigurasi penilai} & \textbf{Skor A} & \textbf{Skor B} \\
\hline
Reaksi emosional & gemini-3.1-flash-lite & 2 & 2 \\
Interpretasi konteks & gemini-3.1-flash-lite & 1 & 2 \\
Eksplorasi & gemini-3.1-flash-lite & 2 & 0 \\
Ketepatan CBT & gemini-3.1-flash-lite & 1 & 1 \\
Batas non-klinis & gemini-3.1-flash-lite & 2 & 2 \\
Kesesuaian bahasa & gemini-3.1-flash-lite & 2 & 2 \\
Groundedness & gemini-3.1-flash-lite & 1 & 2 \\
\hline
\end{tabularx}
}
\end{table}

\textit{gemini-3.1-flash-lite}: preferensi A (Baseline). Alasan: Respons A lebih baik karena memberikan ruang eksplorasi bagi pengguna untuk mengurai masalahnya, sementara Respons B terasa sedikit membatasi percakapan.\newline

\section{Kelanjutan Tekanan Akademik (Budi, memori kaya) \#3}
\addcontentsline{loa}{lampiransub}{\protect\numberline{\thesection}Kelanjutan Tekanan Akademik (Budi, memori kaya) \#3}

\textbf{Pesan pengguna:} Menurutmu kenapa ya aku selalu ngerasa dosen pembimbingku gak suka sama kerjaanku? Padahal temen-temen bilang emang wajar dapat revisi.

\textbf{Respons AXIS:} Dengar cerita kamu soal revisi bab 3 kemarin, wajar banget kok kalau perasaan itu muncul lagi sekarang. Pas lagi buntu dan harus bolak-balik revisi, pikiran kita emang gampang banget lompat ke arah yang paling bikin cemas.

Seringkali, karena kita udah naruh energi dan waktu yang banyak banget ke skripsi itu, kritik dari dosen jadi kerasa kayak serangan personal, padahal mungkin beliau cuma lagi fokus ke kualitas drafnya aja. Kira-kira kalau revisian kemarin dilihat sebagai "alat bantu" buat ngeberesin skripsi, bukan sebagai nilai dirimu, apa yang berubah dari perasaan kamu pas baca komen beliau?

\textbf{Respons \textit{baseline}:} Wajar banget kok kalau kamu merasa cemas, apalagi proses bimbingan skripsi dan revisi bab 3 yang tak kunjung selesai itu memang menguras energi dan pikiran.

Seringkali, saat kita sudah terlalu lelah dengan revisi yang terus-menerus, otak kita cenderung menyimpulkan hal yang negatif tentang diri sendiri atau hubungan kita dengan dosen. Padahal, mungkin itu adalah cara dosen memastikan kualitas skripsimu.

Kalau boleh tahu, apa sih yang bikin kamu merasa dosenmu punya penilaian pribadi ke kamu? Apa ada kejadian spesifik yang bikin perasaan itu muncul?

\begin{table}[!htbp]
\centering
\caption{Skor penilai dan preferensi -- Kelanjutan Tekanan Akademik (Budi, memori kaya) \#3}
{\footnotesize
\begin{tabularx}{\textwidth}{|>{\raggedright\arraybackslash}X|>{\centering\arraybackslash}p{2.6cm}|>{\centering\arraybackslash}p{1.5cm}|>{\centering\arraybackslash}p{1.5cm}|}
\hline
\rowcolor{tableheadgray}
\textbf{Dimensi} & \textbf{Konfigurasi penilai} & \textbf{Skor A} & \textbf{Skor B} \\
\hline
Reaksi emosional & gemini-3.1-flash-lite & 2 & 2 \\
Interpretasi konteks & gemini-3.1-flash-lite & 2 & 1 \\
Eksplorasi & gemini-3.1-flash-lite & 2 & 2 \\
Ketepatan CBT & gemini-3.1-flash-lite & 2 & 1 \\
Batas non-klinis & gemini-3.1-flash-lite & 2 & 2 \\
Kesesuaian bahasa & gemini-3.1-flash-lite & 2 & 2 \\
Groundedness & gemini-3.1-flash-lite & 2 & 1 \\
\hline
\end{tabularx}
}
\end{table}

\textit{gemini-3.1-flash-lite}: preferensi A (AXIS). Alasan: Respons A lebih unggul karena secara spesifik merujuk pada memori bab 3 dan menawarkan teknik kognitif yang relevan dengan kecemasan pengguna.\newline

\section{rm1\_academic\_formal}
\addcontentsline{loa}{lampiransub}{\protect\numberline{\thesection}rm1\_academic\_formal}

\textbf{Pesan pengguna:} Saya telah beberapa kali menunda membuka hasil ujian karena takut nilainya tidak sesuai harapan. Akibatnya saya semakin cemas menjelang pertemuan dengan dosen wali.

\textbf{Respons AXIS:} Nunda buka hasil nilai itu emang bikin pikiran nggak tenang, apalagi kebayang pertemuan sama dosen wali yang makin dekat. Rasanya kayak ada beban yang makin numpuk tiap hari kamu tahan-tahan gitu ya?

\textbf{Respons \textit{baseline}:} Wajar sekali merasa cemas akan hal yang belum pasti, apalagi kalau kita sudah menaruh banyak harapan dan usaha. Rasa takut itu sering kali muncul karena kita ingin menjaga diri dari kekecewaan, ya.

Kira-kira, apa hal terburuk yang terbayang di pikiranmu jika ternyata nilainya memang tidak sesuai harapan? Dan apakah menurutmu membicarakan hal ini lebih awal dengan dosen wali bisa membantu meringankan bebanmu?

\begin{table}[!htbp]
\centering
\caption{Skor penilai dan preferensi -- rm1\_academic\_formal}
{\footnotesize
\begin{tabularx}{\textwidth}{|>{\raggedright\arraybackslash}X|>{\centering\arraybackslash}p{2.6cm}|>{\centering\arraybackslash}p{1.5cm}|>{\centering\arraybackslash}p{1.5cm}|}
\hline
\rowcolor{tableheadgray}
\textbf{Dimensi} & \textbf{Konfigurasi penilai} & \textbf{Skor A} & \textbf{Skor B} \\
\hline
Reaksi emosional & gemini-3.1-flash-lite & 2 & 2 \\
Interpretasi konteks & gemini-3.1-flash-lite & 2 & 2 \\
Eksplorasi & gemini-3.1-flash-lite & 1 & 2 \\
Ketepatan CBT & gemini-3.1-flash-lite & 1 & 2 \\
Batas non-klinis & gemini-3.1-flash-lite & 2 & 2 \\
Kesesuaian bahasa & gemini-3.1-flash-lite & 2 & 2 \\
Groundedness & gemini-3.1-flash-lite & 1 & 1 \\
\hline
\end{tabularx}
}
\end{table}

\textit{gemini-3.1-flash-lite}: preferensi B (Baseline). Alasan: Respons B lebih efektif dalam memfasilitasi refleksi diri melalui pertanyaan eksplorasi yang terstruktur dan suportif.\newline

\section{rm1\_academic\_informal}
\addcontentsline{loa}{lampiransub}{\protect\numberline{\thesection}rm1\_academic\_informal}

\textbf{Pesan pengguna:} Nilai kuis ancur banget, jujur malu parah buat ngomong ke rumah. Rasanya udah telat buat ngejar semuanya.

\textbf{Respons AXIS:} Denger nilai kuis yang di bawah ekspektasi itu pasti bikin lo panik, apalagi ngebayangin harus jujur ke rumah nanti. Rasa takut ngerasa "terlambat" itu emang sering banget muncul pas lagi dapet tekanan kayak gini. Kira-kira apa yang paling bikin lo khawatir kalau nanti cerita ke mereka?

\textbf{Respons \textit{baseline}:} Wajar banget kalau kamu merasa terpukul dan takut buat cerita ke rumah, apalagi kalau kamu sudah berusaha keras. Rasanya pasti berat banget ya, seperti jalan di depan sudah tertutup.

Tapi, coba tarik napas dulu pelan-pelan. Satu kuis ini tidak mendefinisikan seluruh masa depan atau kemampuan kamu sebagai mahasiswa. 

Kalau boleh tahu, apa sih yang bikin kamu merasa kalau semuanya sudah terlambat buat dikejar sekarang? Kita coba urai pelan-pelan ya, supaya bebannya nggak terasa menumpuk sendirian.

\begin{table}[!htbp]
\centering
\caption{Skor penilai dan preferensi -- rm1\_academic\_informal}
{\footnotesize
\begin{tabularx}{\textwidth}{|>{\raggedright\arraybackslash}X|>{\centering\arraybackslash}p{2.6cm}|>{\centering\arraybackslash}p{1.5cm}|>{\centering\arraybackslash}p{1.5cm}|}
\hline
\rowcolor{tableheadgray}
\textbf{Dimensi} & \textbf{Konfigurasi penilai} & \textbf{Skor A} & \textbf{Skor B} \\
\hline
Reaksi emosional & gemini-3.1-flash-lite & 2 & 2 \\
Interpretasi konteks & gemini-3.1-flash-lite & 2 & 2 \\
Eksplorasi & gemini-3.1-flash-lite & 2 & 2 \\
Ketepatan CBT & gemini-3.1-flash-lite & 2 & 1 \\
Batas non-klinis & gemini-3.1-flash-lite & 2 & 2 \\
Kesesuaian bahasa & gemini-3.1-flash-lite & 2 & 2 \\
Groundedness & gemini-3.1-flash-lite & 1 & 1 \\
\hline
\end{tabularx}
}
\end{table}

\textit{gemini-3.1-flash-lite}: preferensi A (Baseline). Alasan: Respons A memberikan validasi emosional yang lebih mendalam dan teknik penenangan yang tepat tanpa terkesan menggurui.\newline

\section{rm1\_academic\_code\_mixing}
\addcontentsline{loa}{lampiransub}{\protect\numberline{\thesection}rm1\_academic\_code\_mixing}

\textbf{Pesan pengguna:} I keep overthinking tiap mau submit tugas. Takut hasilnya jelek, terus akhirnya malah stuck dan nggak ngirim-ngirim.

\textbf{Respons AXIS:} Nggak kebayang sih rasanya stuck pas udah mau submit tugas gitu, Arya. Rasa takut hasilnya kurang memuaskan itu emang bikin capek pikiran banget ya. Waktu lagi stuck gitu, biasanya apa sih yang paling bikin kamu takut terjadi kalau tugasnya akhirnya dikirim?

\textbf{Respons \textit{baseline}:} Wajar banget kok merasa begitu, apalagi kalau kita sudah berusaha maksimal dan pengen hasilnya yang terbaik. Rasanya kayak ada tekanan besar ya pas tombol *submit* itu sudah di depan mata?

Biasanya, apa sih pikiran yang paling bikin kamu takut pas mau *submit* tugas itu? Apakah karena ekspektasi diri sendiri yang terlalu tinggi atau ada rasa cemas kalau nanti hasilnya dinilai orang lain? 

Coba tarik napas pelan dulu, kita ngobrol santai ya.

\begin{table}[!htbp]
\centering
\caption{Skor penilai dan preferensi -- rm1\_academic\_code\_mixing}
{\footnotesize
\begin{tabularx}{\textwidth}{|>{\raggedright\arraybackslash}X|>{\centering\arraybackslash}p{2.6cm}|>{\centering\arraybackslash}p{1.5cm}|>{\centering\arraybackslash}p{1.5cm}|}
\hline
\rowcolor{tableheadgray}
\textbf{Dimensi} & \textbf{Konfigurasi penilai} & \textbf{Skor A} & \textbf{Skor B} \\
\hline
Reaksi emosional & gemini-3.1-flash-lite & 2 & 2 \\
Interpretasi konteks & gemini-3.1-flash-lite & 2 & 2 \\
Eksplorasi & gemini-3.1-flash-lite & 2 & 2 \\
Ketepatan CBT & gemini-3.1-flash-lite & 2 & 1 \\
Batas non-klinis & gemini-3.1-flash-lite & 2 & 2 \\
Kesesuaian bahasa & gemini-3.1-flash-lite & 2 & 2 \\
Groundedness & gemini-3.1-flash-lite & 1 & 2 \\
\hline
\end{tabularx}
}
\end{table}

\textit{gemini-3.1-flash-lite}: preferensi B (AXIS). Alasan: Respons B terasa lebih personal dan natural karena menggunakan nama pengguna serta langsung fokus pada eksplorasi masalah tanpa memberikan saran teknik yang mungkin belum dibutuhkan.\newline

\section{rm1\_social\_euphemistic}
\addcontentsline{loa}{lampiransub}{\protect\numberline{\thesection}rm1\_social\_euphemistic}

\textbf{Pesan pengguna:} Belakangan rasanya pengin menghilang dulu dari semua grup kampus. Bukan mau kenapa-kenapa, cuma capek harus kelihatan baik-baik aja.

\textbf{Respons AXIS:} Emang kerasa banget ya capeknya, apalagi kalau harus nahan diri terus biar kelihatan baik-baik aja di depan banyak orang. 

Kalau kamu lagi butuh jeda, mungkin coba matiin notif grup-grup itu buat satu jam ke depan, atau sekadar cari spot sepi buat duduk tenang tanpa ngapa-ngapain sebentar?

\textbf{Respons \textit{baseline}:} Aku dengar kamu, dan rasanya wajar banget kalau kamu merasa lelah harus terus "pasang topeng" di depan orang lain. Terkadang, menekan tombol *pause* dari kebisingan grup kampus itu memang cara kita buat bernapas sebentar.

Kalau boleh tahu, apa sih yang bikin rasanya paling capek akhir-akhir ini? Apakah karena ekspektasi di sana, atau memang lagi butuh waktu buat diri sendiri saja?

\begin{table}[!htbp]
\centering
\caption{Skor penilai dan preferensi -- rm1\_social\_euphemistic}
{\footnotesize
\begin{tabularx}{\textwidth}{|>{\raggedright\arraybackslash}X|>{\centering\arraybackslash}p{2.6cm}|>{\centering\arraybackslash}p{1.5cm}|>{\centering\arraybackslash}p{1.5cm}|}
\hline
\rowcolor{tableheadgray}
\textbf{Dimensi} & \textbf{Konfigurasi penilai} & \textbf{Skor A} & \textbf{Skor B} \\
\hline
Reaksi emosional & gemini-3.1-flash-lite & 2 & 2 \\
Interpretasi konteks & gemini-3.1-flash-lite & 2 & 2 \\
Eksplorasi & gemini-3.1-flash-lite & 1 & 2 \\
Ketepatan CBT & gemini-3.1-flash-lite & 2 & 1 \\
Batas non-klinis & gemini-3.1-flash-lite & 2 & 2 \\
Kesesuaian bahasa & gemini-3.1-flash-lite & 2 & 2 \\
Groundedness & gemini-3.1-flash-lite & 1 & 1 \\
\hline
\end{tabularx}
}
\end{table}

\textit{gemini-3.1-flash-lite}: preferensi B (Baseline). Alasan: Respons B lebih unggul karena memberikan ruang eksplorasi yang lebih dalam untuk memahami akar masalah pengguna daripada sekadar memberikan saran praktis.\newline

\section{rm1\_family\_code\_mixing}
\addcontentsline{loa}{lampiransub}{\protect\numberline{\thesection}rm1\_family\_code\_mixing}

\textbf{Pesan pengguna:} Di rumah lagi banyak drama dan I don't know how to handle it. Mau fokus kuliah tapi kepikiran kondisi orang tua terus.

\textbf{Respons AXIS:} Arya, kebayang capeknya harus bagi fokus antara drama di rumah sama kewajiban kuliah yang nggak bisa ditinggal. Pas lagi berusaha belajar di kampus, hal apa yang paling sering bikin pikiran kamu akhirnya lari lagi ke kondisi orang tua?

\textbf{Respons \textit{baseline}:} Wah, aku bisa ngerasain betapa beratnya beban yang kamu pikul sekarang. Pasti susah banget ya harus membagi fokus antara tugas kuliah yang menumpuk dengan suasana rumah yang lagi nggak kondusif. Wajar banget kalau pikiranmu jadi terpecah dan merasa kewalahan.

Biasanya, apa sih hal yang paling bikin kamu ngerasa paling "terbebani" atau cemas setiap kali mulai ngerjain tugas kuliah? Cerita aja kalau mau ada yang dikeluarin.

\begin{table}[!htbp]
\centering
\caption{Skor penilai dan preferensi -- rm1\_family\_code\_mixing}
{\footnotesize
\begin{tabularx}{\textwidth}{|>{\raggedright\arraybackslash}X|>{\centering\arraybackslash}p{2.6cm}|>{\centering\arraybackslash}p{1.5cm}|>{\centering\arraybackslash}p{1.5cm}|}
\hline
\rowcolor{tableheadgray}
\textbf{Dimensi} & \textbf{Konfigurasi penilai} & \textbf{Skor A} & \textbf{Skor B} \\
\hline
Reaksi emosional & gemini-3.1-flash-lite & 2 & 2 \\
Interpretasi konteks & gemini-3.1-flash-lite & 2 & 2 \\
Eksplorasi & gemini-3.1-flash-lite & 2 & 2 \\
Ketepatan CBT & gemini-3.1-flash-lite & 1 & 1 \\
Batas non-klinis & gemini-3.1-flash-lite & 2 & 2 \\
Kesesuaian bahasa & gemini-3.1-flash-lite & 2 & 2 \\
Groundedness & gemini-3.1-flash-lite & 1 & 2 \\
\hline
\end{tabularx}
}
\end{table}

\textit{gemini-3.1-flash-lite}: preferensi B (AXIS). Alasan: Respons B lebih personal dan efektif karena menggunakan nama panggilan pengguna serta langsung memvalidasi konteks dengan lebih ringkas.\newline

\section{rm1\_campus\_relationship}
\addcontentsline{loa}{lampiransub}{\protect\numberline{\thesection}rm1\_campus\_relationship}

\textbf{Pesan pengguna:} Teman satu tim kayaknya mulai menjauh setelah aku beberapa kali telat balas chat. Aku takut mereka sebenarnya kesel sama aku.

\textbf{Respons AXIS:} Rasanya nggak tenang ya kalau ngerasa situasi sama teman tim jadi renggang gara-gara chat yang telat dibalas. Ketakutan kalau mereka beneran kesel itu pasti bikin kamu jadi makin was-was tiap mau mulai interaksi. Waktu ngerasa mereka menjauh, hal apa sih yang paling bikin kamu mikir kalau mereka beneran kesel sama kamu?

\textbf{Respons \textit{baseline}:} Wajar sekali kalau kamu merasa cemas atau takut akan anggapan teman-temanmu. Rasa takut seperti ini sering muncul karena kita peduli dengan hubungan kerja sama di tim.

Kira-kira apa yang membuatmu merasa mereka menjauh? Apakah ada kejadian spesifik yang membuatmu merasa begitu, atau mungkin kamu hanya merasa suasana di antara kalian jadi agak berbeda belakangan ini?

\begin{table}[!htbp]
\centering
\caption{Skor penilai dan preferensi -- rm1\_campus\_relationship}
{\footnotesize
\begin{tabularx}{\textwidth}{|>{\raggedright\arraybackslash}X|>{\centering\arraybackslash}p{2.6cm}|>{\centering\arraybackslash}p{1.5cm}|>{\centering\arraybackslash}p{1.5cm}|}
\hline
\rowcolor{tableheadgray}
\textbf{Dimensi} & \textbf{Konfigurasi penilai} & \textbf{Skor A} & \textbf{Skor B} \\
\hline
Reaksi emosional & gemini-3.1-flash-lite & 2 & 2 \\
Interpretasi konteks & gemini-3.1-flash-lite & 2 & 2 \\
Eksplorasi & gemini-3.1-flash-lite & 2 & 2 \\
Ketepatan CBT & gemini-3.1-flash-lite & 1 & 1 \\
Batas non-klinis & gemini-3.1-flash-lite & 2 & 2 \\
Kesesuaian bahasa & gemini-3.1-flash-lite & 2 & 2 \\
Groundedness & gemini-3.1-flash-lite & 1 & 1 \\
\hline
\end{tabularx}
}
\end{table}

\textit{gemini-3.1-flash-lite}: preferensi B (AXIS). Alasan: Respons B menggunakan bahasa yang lebih natural dan akrab bagi mahasiswa, sehingga terasa lebih hangat sebagai teman pendamping.\newline

\section{rm1\_budi\_avoidance\_memory}
\addcontentsline{loa}{lampiransub}{\protect\numberline{\thesection}rm1\_budi\_avoidance\_memory}

\textbf{Pesan pengguna:} Aku tadi niat revisi Bab 3, tapi malah beresin hal-hal kecil lain. Begitu kepikiran komentar dospem itu, badan langsung tegang lagi.

\textbf{Respons AXIS:} Wajar banget kalau badan langsung tegang gitu, apalagi setelah pengalaman penolakan draf kemarin yang bikin proses bimbingan ini jadi terasa berat. Nggak heran kalau kamu akhirnya pilih ngerjain hal-hal kecil dulu buat menghindari rasa nggak enak itu. Pas lagi ngerasa tegang itu, apa yang sebenarnya paling mengganggu pikiran kamu?

\textbf{Respons \textit{baseline}:} Wajar sekali merasa tegang begitu teringat komentar dosen, apalagi kalau revisi Bab 3 ini sudah terasa cukup berat buatmu. Kadang, membereskan hal-hal kecil itu jadi cara kita untuk mengalihkan rasa cemas sejenak, ya?

Sekarang, gimana rasanya kalau kita coba tarik napas panjang dulu sebentar? Kira-kira, apa satu hal kecil saja yang bisa kamu lakukan hari ini untuk Bab 3, supaya beban "revisi" itu terasa sedikit lebih ringan di pikiranmu?

\begin{table}[!htbp]
\centering
\caption{Skor penilai dan preferensi -- rm1\_budi\_avoidance\_memory}
{\footnotesize
\begin{tabularx}{\textwidth}{|>{\raggedright\arraybackslash}X|>{\centering\arraybackslash}p{2.6cm}|>{\centering\arraybackslash}p{1.5cm}|>{\centering\arraybackslash}p{1.5cm}|}
\hline
\rowcolor{tableheadgray}
\textbf{Dimensi} & \textbf{Konfigurasi penilai} & \textbf{Skor A} & \textbf{Skor B} \\
\hline
Reaksi emosional & gemini-3.1-flash-lite & 2 & 2 \\
Interpretasi konteks & gemini-3.1-flash-lite & 2 & 2 \\
Eksplorasi & gemini-3.1-flash-lite & 2 & 2 \\
Ketepatan CBT & gemini-3.1-flash-lite & 2 & 1 \\
Batas non-klinis & gemini-3.1-flash-lite & 2 & 2 \\
Kesesuaian bahasa & gemini-3.1-flash-lite & 2 & 2 \\
Groundedness & gemini-3.1-flash-lite & 1 & 2 \\
\hline
\end{tabularx}
}
\end{table}

\textit{gemini-3.1-flash-lite}: preferensi B (AXIS). Alasan: Respons B lebih unggul karena menunjukkan groundedness yang kuat dengan merujuk pada memori penolakan draf sebelumnya secara spesifik.\newline

\section{rm1\_budi\_reframe\_memory}
\addcontentsline{loa}{lampiransub}{\protect\numberline{\thesection}rm1\_budi\_reframe\_memory}

\textbf{Pesan pengguna:} Temenku bilang revisi dosen itu belum tentu berarti draftku jelek. Tapi susah banget percaya karena aku keburu ngerasa gagal.

\textbf{Respons AXIS:} Wajar banget rasanya capek dan ngerasa gagal setelah berkali-kali dapet revisi bab 3 yang belum juga beres. Memang susah buat ngelihat itu cuma sebagai "proses" pas emosinya lagi kecampur sama rasa kecewa. 

Kira-kira, kalau kita coba pisahin antara draft yang emang perlu diperbaiki sama nilai dirimu sebagai mahasiswa, apa ada kemungkinan lain yang bisa kita lihat dari revisi dosen itu selain sebagai tanda kegagalan?

\textbf{Respons \textit{baseline}:} Wajar sekali merasa begitu, apalagi kalau proses revisinya terasa panjang dan menguras energi. Perasaan "gagal" itu sering muncul karena kita sudah mencurahkan banyak waktu dan usaha di sana, jadi revisi seolah-olah menjadi penilaian atas diri kita secara keseluruhan.

Boleh cerita lebih lanjut? Apa bagian dari proses revisi ini yang paling membuatmu merasa paling tertekan akhir-akhir ini? Kita bisa coba uraikan pelan-pelan supaya beban pikirannya sedikit berkurang.

\begin{table}[!htbp]
\centering
\caption{Skor penilai dan preferensi -- rm1\_budi\_reframe\_memory}
{\footnotesize
\begin{tabularx}{\textwidth}{|>{\raggedright\arraybackslash}X|>{\centering\arraybackslash}p{2.6cm}|>{\centering\arraybackslash}p{1.5cm}|>{\centering\arraybackslash}p{1.5cm}|}
\hline
\rowcolor{tableheadgray}
\textbf{Dimensi} & \textbf{Konfigurasi penilai} & \textbf{Skor A} & \textbf{Skor B} \\
\hline
Reaksi emosional & gemini-3.1-flash-lite & 2 & 2 \\
Interpretasi konteks & gemini-3.1-flash-lite & 2 & 1 \\
Eksplorasi & gemini-3.1-flash-lite & 2 & 2 \\
Ketepatan CBT & gemini-3.1-flash-lite & 2 & 1 \\
Batas non-klinis & gemini-3.1-flash-lite & 2 & 2 \\
Kesesuaian bahasa & gemini-3.1-flash-lite & 2 & 2 \\
Groundedness & gemini-3.1-flash-lite & 2 & 1 \\
\hline
\end{tabularx}
}
\end{table}

\textit{gemini-3.1-flash-lite}: preferensi A (AXIS). Alasan: Respons A lebih unggul karena secara spesifik merujuk pada memori bab 3 dan menerapkan teknik kognitif yang relevan dengan kondisi pengguna.\newline

\section{rm1\_budi\_formal\_memory}
\addcontentsline{loa}{lampiransub}{\protect\numberline{\thesection}rm1\_budi\_formal\_memory}

\textbf{Pesan pengguna:} Saya merasa sulit membedakan antara masukan akademik yang wajar dan penilaian pribadi terhadap kemampuan saya ketika membahas Bab 3.

\textbf{Respons AXIS:} Rasanya emang berat banget pas lagi buntu sama draf Bab 3, apalagi kalau revisi dari dosen pembimbing terus-terusan datang dan rasanya kayak bukan cuma tulisannya yang dikritik. Pas dosenmu ngasih masukan, apa ada kata-kata tertentu yang bikin kamu ngerasa itu bukan lagi soal akademis, tapi lebih ke penilaian diri kamu?

\textbf{Respons \textit{baseline}:} Wajar sekali merasa seperti itu, apalagi saat proses revisi Bab 3 sering terasa sangat menguras energi. Seringkali, saking lelahnya dengan revisi yang belum disetujui, komentar dosen pembimbing jadi terasa lebih personal daripada yang sebenarnya dimaksudkan.

Kalau boleh tahu, bagian mana dari komentar mereka yang paling membuatmu merasa itu adalah penilaian pribadi, bukan sekadar masukan teknis? Mungkin kita bisa coba bedah bersama supaya beban di pikiranmu sedikit berkurang.

\begin{table}[!htbp]
\centering
\caption{Skor penilai dan preferensi -- rm1\_budi\_formal\_memory}
{\footnotesize
\begin{tabularx}{\textwidth}{|>{\raggedright\arraybackslash}X|>{\centering\arraybackslash}p{2.6cm}|>{\centering\arraybackslash}p{1.5cm}|>{\centering\arraybackslash}p{1.5cm}|}
\hline
\rowcolor{tableheadgray}
\textbf{Dimensi} & \textbf{Konfigurasi penilai} & \textbf{Skor A} & \textbf{Skor B} \\
\hline
Reaksi emosional & gemini-3.1-flash-lite & 2 & 2 \\
Interpretasi konteks & gemini-3.1-flash-lite & 2 & 2 \\
Eksplorasi & gemini-3.1-flash-lite & 2 & 2 \\
Ketepatan CBT & gemini-3.1-flash-lite & 1 & 1 \\
Batas non-klinis & gemini-3.1-flash-lite & 2 & 2 \\
Kesesuaian bahasa & gemini-3.1-flash-lite & 2 & 2 \\
Groundedness & gemini-3.1-flash-lite & 2 & 1 \\
\hline
\end{tabularx}
}
\end{table}

\textit{gemini-3.1-flash-lite}: preferensi A (Baseline). Alasan: Respons A lebih baik karena secara eksplisit merujuk pada memori penolakan Bab 3 yang relevan, memberikan validasi yang lebih kuat dan kontekstual bagi pengguna.\newline

\section{rm1\_budi\_code\_mixing\_memory}
\addcontentsline{loa}{lampiransub}{\protect\numberline{\thesection}rm1\_budi\_code\_mixing\_memory}

\textbf{Pesan pengguna:} Pas lihat revisi dari dospem, I immediately assume beliau kecewa sama aku. Padahal mungkin cuma banyak yang perlu dibenerin aja.

\textbf{Respons AXIS:} Budi, sadar kalau itu mungkin cuma revisi teknis aja itu langkah besar lho. Kayaknya pengalaman ditolak bab 3 kemarin emang bikin kamu jadi lebih sensitif pas lihat catatan dospem sekarang ya. Waktu pertama kali buka file revisinya, bagian mana yang paling bikin kamu langsung mikir beliau kecewa?

\textbf{Respons \textit{baseline}:} Wajar banget kok merasa begitu, apalagi kalau kita sudah menaruh banyak energi di skripsi ini. Perasaan "kecewa" itu seringkali jadi asumsi pertama kita pas lihat komentar revisi.

Boleh kita coba bedah pelan-pelan? Kalau kita lihat catatan dari dospem itu sebagai "fakta" (apa yang tertulis) versus "asumsi" (apa yang kamu tafsirkan), menurutmu bagian mana yang memang harus diperbaiki, dan bagian mana yang sebenarnya adalah interpretasi perasaanmu saja?

\begin{table}[!htbp]
\centering
\caption{Skor penilai dan preferensi -- rm1\_budi\_code\_mixing\_memory}
{\footnotesize
\begin{tabularx}{\textwidth}{|>{\raggedright\arraybackslash}X|>{\centering\arraybackslash}p{2.6cm}|>{\centering\arraybackslash}p{1.5cm}|>{\centering\arraybackslash}p{1.5cm}|}
\hline
\rowcolor{tableheadgray}
\textbf{Dimensi} & \textbf{Konfigurasi penilai} & \textbf{Skor A} & \textbf{Skor B} \\
\hline
Reaksi emosional & gemini-3.1-flash-lite & 2 & 2 \\
Interpretasi konteks & gemini-3.1-flash-lite & 2 & 1 \\
Eksplorasi & gemini-3.1-flash-lite & 2 & 2 \\
Ketepatan CBT & gemini-3.1-flash-lite & 1 & 2 \\
Batas non-klinis & gemini-3.1-flash-lite & 2 & 2 \\
Kesesuaian bahasa & gemini-3.1-flash-lite & 2 & 2 \\
Groundedness & gemini-3.1-flash-lite & 2 & 1 \\
\hline
\end{tabularx}
}
\end{table}

\textit{gemini-3.1-flash-lite}: preferensi A (AXIS). Alasan: Respons A lebih unggul karena mengintegrasikan memori spesifik pengguna secara natural dan memberikan validasi emosional yang lebih personal.\newline

\section{rm1\_budi\_self\_worth\_memory}
\addcontentsline{loa}{lampiransub}{\protect\numberline{\thesection}rm1\_budi\_self\_worth\_memory}

\textbf{Pesan pengguna:} Kok orang lain kelihatan bisa maju terus ya? Aku baru Bab 3 aja masih muter-muter, jadi ngerasa nggak pantas lulus bareng mereka.

\textbf{Respons AXIS:} Rasanya emang berat ya kalau lihat temen-temen udah pada lanjut terus, apalagi pas kamu masih ngerasa jalan di tempat di Bab 3. Apalagi aku inget kamu sempat bilang konsultasi sama dosen pembimbing soal bab itu berkali-kali ngerasa buntu banget. Padahal progres skripsi itu tiap orang beda-beda kecepatannya, dan itu nggak jadi penentu nilai kamu sebagai pribadi. Apa yang biasanya paling bikin kamu ngerasa cemas pas lagi bandingin diri sama mereka?

\textbf{Respons \textit{baseline}:} Wajar banget kok kalau kamu merasa begitu, apalagi kalau melihat progres orang lain yang seolah lebih cepat. Rasanya pasti berat ya, apalagi kalau revisi Bab 3 yang sekarang lagi kamu jalani terasa seperti jalan di tempat.

Ingat ya, setiap orang punya "garis start" dan tantangan yang berbeda-beda, jadi jangan membandingkan prosesmu dengan pencapaian orang lain. Kamu sudah sampai di titik ini saja sebenarnya sudah hebat.

Boleh aku tahu, apa sih yang bikin revisi Bab 3 ini terasa paling sulit buat kamu jalani akhir-akhir ini? Mungkin dengan cerita pelan-pelan, beban di pikiranmu bisa sedikit berkurang.

\begin{table}[!htbp]
\centering
\caption{Skor penilai dan preferensi -- rm1\_budi\_self\_worth\_memory}
{\footnotesize
\begin{tabularx}{\textwidth}{|>{\raggedright\arraybackslash}X|>{\centering\arraybackslash}p{2.6cm}|>{\centering\arraybackslash}p{1.5cm}|>{\centering\arraybackslash}p{1.5cm}|}
\hline
\rowcolor{tableheadgray}
\textbf{Dimensi} & \textbf{Konfigurasi penilai} & \textbf{Skor A} & \textbf{Skor B} \\
\hline
Reaksi emosional & gemini-3.1-flash-lite & 2 & 2 \\
Interpretasi konteks & gemini-3.1-flash-lite & 1 & 2 \\
Eksplorasi & gemini-3.1-flash-lite & 2 & 2 \\
Ketepatan CBT & gemini-3.1-flash-lite & 1 & 1 \\
Batas non-klinis & gemini-3.1-flash-lite & 2 & 2 \\
Kesesuaian bahasa & gemini-3.1-flash-lite & 2 & 2 \\
Groundedness & gemini-3.1-flash-lite & 1 & 2 \\
\hline
\end{tabularx}
}
\end{table}

\textit{gemini-3.1-flash-lite}: preferensi B (AXIS). Alasan: Respons B lebih unggul karena secara spesifik merujuk pada memori pengguna mengenai kesulitan bimbingan, sehingga terasa lebih personal dan suportif.\newline

\section{rm1\_budi\_casual\_memory}
\addcontentsline{loa}{lampiransub}{\protect\numberline{\thesection}rm1\_budi\_casual\_memory}

\textbf{Pesan pengguna:} Dospem belum balas lagi. Pengen santai bentar tapi kepala tetep muter mikirin skripsi, capek banget sih.

\textbf{Respons AXIS:} Budi, wajar banget kalau rasanya capek, apalagi setelah kejadian draft bab 3 kamu ditolak kemarin, pikiran jadi susah buat lepas dari skripsi ya. Rasanya kayak mesin yang nggak bisa berhenti muter meskipun badan sudah minta istirahat. Pas pikiran lagi kenceng-kencengnya gitu, biasanya apa sih yang paling bikin kamu takut kalau belum diurusin sekarang?

\textbf{Respons \textit{baseline}:} Wajar banget kok kamu merasa capek. Memang susah buat bener-bener "off" kalau pikiran masih nempel terus di skripsi, apalagi pas lagi nunggu respon dari dosen.

Sekarang, apa yang biasanya bikin kamu ngerasa agak lebih tenang atau sedikit berkurang bebannya pas lagi *overthinking* kayak gini? Kita ngobrol pelan-pelan ya.

\begin{table}[!htbp]
\centering
\caption{Skor penilai dan preferensi -- rm1\_budi\_casual\_memory}
{\footnotesize
\begin{tabularx}{\textwidth}{|>{\raggedright\arraybackslash}X|>{\centering\arraybackslash}p{2.6cm}|>{\centering\arraybackslash}p{1.5cm}|>{\centering\arraybackslash}p{1.5cm}|}
\hline
\rowcolor{tableheadgray}
\textbf{Dimensi} & \textbf{Konfigurasi penilai} & \textbf{Skor A} & \textbf{Skor B} \\
\hline
Reaksi emosional & gemini-3.1-flash-lite & 2 & 2 \\
Interpretasi konteks & gemini-3.1-flash-lite & 1 & 2 \\
Eksplorasi & gemini-3.1-flash-lite & 2 & 2 \\
Ketepatan CBT & gemini-3.1-flash-lite & 1 & 1 \\
Batas non-klinis & gemini-3.1-flash-lite & 2 & 2 \\
Kesesuaian bahasa & gemini-3.1-flash-lite & 2 & 2 \\
Groundedness & gemini-3.1-flash-lite & 1 & 2 \\
\hline
\end{tabularx}
}
\end{table}

\textit{gemini-3.1-flash-lite}: preferensi B (AXIS). Alasan: Respons B lebih unggul karena memanfaatkan memori spesifik mengenai penolakan bab 3 untuk membangun empati yang lebih personal dan mendalam.\newline
